% source: https://github.com/pfradin/luadraw/discussions/240#discussioncomment-16649133 (pfradin, block 1)
\documentclass[12pt]{standalone}
\usepackage[b]{esvect}
\usepackage[svgnames]{xcolor}
\usepackage[3d]{luadraw}
%\author{Huỳnh Văn Thơ}
\begin{document}

\begin{luadraw}{name=fill}
local g = graph:new{window = {-2, 6, -4, 12}, size ={14,14}}
g:Linejoin("round")
local i = cpx.I 

g:Dgrid({-1-4*i, 5+3*i}, {gridstyle="solid"})
g:Daxes({0,1,1}, {grid=true, arrows="-Stealth", legend={"$x$","$y$"}})
--
local f = function(x) return x*x - 2*x - 2 end
local h = function(x) return x*x + 2 end
local k = function(x) return -x*x+5*x + 2 end
--
g:Dcartesian(h, {x={-1,3}, draw_options="blue, line width=0.8pt"})
g:Dcartesian(f, {x={-1,4.2}, draw_options="red, line width=0.8pt"})
g:Dcartesian(k, {x={-1,4.2}, draw_options="green, line width=0.8pt"})
--
local D1 = domain3(f,k,-2,6) -- domain between Cf and Ck
local C = cartesian(h,-2,6)[1] -- Ch curve (first component)
local A, B = C[1], C[#C] -- first and last point of C
table.insert(C,1, -2+i*A.im); table.insert(C,1, -2-4*i)
table.insert(C,6+i*B.im); table.insert(C,6-4*i)
table.insert(C,"l")
-- now C is a path, it is the contour of the part of the plane located under the curve of h, we use it to clip the domain between the other two curves.
g:Beginclip(C)
    g:Dpolyline(D1,"draw=none,fill=pink,fill opacity=0.6")
g:Endclip()
g:Show()
\end{luadraw}
\end{document}
